% ══════════════════════════════════════════════════════════════════════════════
%  EXPERIMENT 2 — SURFACE AREA AND DISSOLVING                    (Lesson 6, Term 2 Week 2)
% ══════════════════════════════════════════════════════════════════════════════
\experimentdiv{Surface Area and Dissolving}{Timing how fast sugar dissolves \textbullet{} Repeat trials and averages \textbullet{} Testing a hypothesis}{Lesson 6 \textbullet\ Term 2, Week 2 \textbullet\ Lab}\label{sec:2}
\lessoncodes{VC2S8U05, U06 \textbullet{} VC2S8I01--I06}

\vocabhead
\begin{ctab}
\begin{tabular}{|SW{4.2cm}|W{11.6cm}|}
  \hline\hd & \thead{Word} & \thead{Definition (copy from the slide)} \\ \hline
  & \scaf{\textbf{Solute}} & \scaf{The substance that dissolves, such as sugar.} \\ \hline
  & \textbf{Solvent} & \\ \hline
  & \textbf{Surface area} & \\ \hline
  & \textbf{Rate} & \\ \hline
  & \textbf{Average} & \\ \hline
\end{tabular}
\end{ctab}

\begin{lipanel}
\inquiryli{Write a hypothesis, time repeat trials and calculate an average.}
\understandli{Explain, using the particle model, why smaller pieces of solute dissolve faster.}
\sctitle
\begin{itemize}[leftmargin=1.4em, itemsep=0pt, topsep=1pt]
  \item I can write a hypothesis: \emph{If the [IV] is \dots, the [DV] will \dots}
  \item I can time three trials and calculate their average.
  \item I can describe the trend in my results and say whether it supports my hypothesis.
\end{itemize}
\end{lipanel}

\begin{skillsfocus}
\focus{QP1}{Write a testable hypothesis.}
\focus{PC2}{Name the independent, dependent and controlled variables.}
\focus{PC8}{Time each trial with a digital stopwatch.}
\focus{PA1}{Record times in a results table.}
\focus{PA3}{Calculate the average of three trials.}
\focus{PA5}{Describe the trend.}
\focus{EV1}{Compare the results with the hypothesis.}
\focus{EV3}{Find sources of error.}
\end{skillsfocus}

% ─── TECHNIQUE ────────────────────────────────────────────────────────────────
\exppart{Technique spotlight}{Blue}

\begin{technique}{Writing a hypothesis}{QP1}
A \textbf{hypothesis} is a prediction your experiment can really test. It names the independent variable (IV) and the dependent variable (DV), and says which way the DV will change:
\begin{center}\emph{If the} [IV] \emph{is} [increased / decreased], \emph{the} [DV] \emph{will} [increase / decrease].\end{center}
A reason is optional but strengthens it: add \emph{\dots because} [the science reason].
\end{technique}

\begin{workedex}{A hypothesis about temperature}
\emph{If the temperature of the water is increased, the time for the sugar to dissolve will decrease.}

\smallskip
With a reason: \emph{\dots because the water particles move faster and hit the sugar more often.}
\end{workedex}

\begin{wedo}{A hypothesis about stirring}
Write a hypothesis for an experiment that tests whether stirring changes how fast salt dissolves.
\answerbox{3.5cm}
\end{wedo}

\begin{technique}{Repeat trials and averages}{PA3}
One timing can be wrong by chance: you might start the stopwatch late. Do each test \textbf{three times} and calculate the \textbf{average}:
\begin{center}average $=$ (trial 1 $+$ trial 2 $+$ trial 3) $\div$ 3\end{center}
If one trial is very different from the other two, it is an \textbf{outlier}. Circle it, leave it out, and average the other two.
\end{technique}

\begin{workedex}{Calculating an average}
Times: 42\,s, 45\,s, 39\,s.\qquad Average $= (42 + 45 + 39) \div 3 = 126 \div 3 = 42$\,s.
\end{workedex}

\begin{youdobox}[Averages]
\task Calculate the average of each set of times. Show your working. \skr{PA3}
\pt{a}55\,s, 61\,s, 58\,s
\workbox{2cm}
\pt{b}20\,s, 23\,s, 47\,s (look for an outlier first)
\workbox{2.5cm}
\end{youdobox}

% ─── PLAN AND SAFETY / METHOD ────────────────────────────────────────────────
\exppart{Plan, safety and method}{Teal}

\textbf{Aim:} To investigate the effect of changing the size of the sugar pieces on the time sugar takes to dissolve.

\task Name the variables in this experiment. \skr{PC2}
\begin{ctab}
\begin{tabular}{|SW{4.4cm}|W{10.8cm}|}
  \hline\hd & \thead{Variable} & \thead{In this experiment} \\ \hline
  & \scaf{Independent} & \scaf{Size of the sugar pieces (cube, granulated, caster)} \\ \hline
  & Dependent & \\ \hline
  & Controlled & \\ \hline
  & Controlled & \\ \hline
\end{tabular}
\end{ctab}

\needspace{8\baselineskip}
\task Write your hypothesis. \skr{QP1}
\answerbox{2.4cm}

\subsubsection*{Equipment}
3 $\times$ 250\,mL beakers \textbullet{} 100\,mL measuring cylinder \textbullet{} electronic balance \textbullet{} stirring rod \textbullet{} digital stopwatch (or phone) \textbullet{} sugar cubes, granulated sugar and caster sugar \textbullet{} room-temperature water

\begin{infobox}{Safety}
Wear safety glasses. Do not eat the sugar: lab equipment is not clean enough for food. Keep the phone away from the water.
\end{infobox}

\subsubsection*{Method}
\begin{enumerate}[itemsep=2pt]
  \item Measure 100\,mL of water into a beaker.
  \item Weigh out 4\,g of sugar: one cube, or 4\,g of granulated or caster sugar.
  \item Tip the sugar into the water and start the stopwatch at the same moment.
  \item Stir at a steady rate, about one stir per second. Stop the stopwatch when no sugar can be seen.
  \item Record the time in Table 2.1, then rinse the beaker.
  \item Repeat for the other two sizes, then do the whole experiment twice more.
\end{enumerate}

% ─── RESULTS ─────────────────────────────────────────────────────────────────
\exppart{Results}{Amber}

\subsubsection*{Table 2.1: Time for 4\,g of sugar to dissolve in 100\,mL of water \skr{PA1}}
\begin{ctab}
\begin{tabular}{|SW{3.6cm}|C{2.3cm}|C{2.3cm}|C{2.3cm}|C{3.0cm}|}
  \hline\hd & \thead{Size of pieces} & \thead{Trial 1 (s)} & \thead{Trial 2 (s)} & \thead{Trial 3 (s)} & \thead{Average (s)} \\ \hline
  & \scaf{Sugar cube} & \scaf{150} & \scaf{162} & \scaf{141} & \scaf{151} \\ \hline
  & Sugar cube & & & & \\ \hline
  & Granulated sugar & & & & \\ \hline
  & Caster sugar & & & & \\ \hline
\end{tabular}
\end{ctab}
{\small\color{DarkGrey}The red row is an example. Record your own results in the rows below it.}

\task Show your working for one of the averages. \skr{PA3}
\workbox{3cm}

\needspace{16\baselineskip}
\task Draw a column graph of the average times. Put the size of pieces on the horizontal axis. Label both axes and give the graph a title. \skr{PA2}
\graphgrid{8cm}

% ─── ANALYSIS ────────────────────────────────────────────────────────────────
\exppart{Analysis and evaluation}{Coral}

\task Describe the trend in your results. Use numbers from your table. \skr{PA5}
\answerbox{4cm}

\needspace{10\baselineskip}
\task Do your results support your hypothesis? State the trend, back it with numbers from Table 2.1, and say \emph{supported} or \emph{not supported} (never \emph{proven}). \skr{EV1}
\answerbox{4cm}

\needspace{12\baselineskip}
\task Use the particle model to explain why caster sugar dissolves faster than a sugar cube. Use the words \emph{surface area}, \emph{solvent particles} and \emph{solute particles}.
\answerbox{5cm}

\needspace{12\baselineskip}
\task Give two sources of error in this experiment, and say how each could be reduced. \skr{EV3}
\begin{ctab}
\begin{tabular}{|SW{7.2cm}|W{7.8cm}|}
  \hline\hd & \thead{Source of error} & \thead{How to reduce it} \\ \hline
  & \scaf{Deciding when the last sugar has gone} & \scaf{Have the same person judge every trial, looking against a white card} \\ \hline
  & & \\ \hline
  & & \\ \hline
\end{tabular}
\end{ctab}

\begin{yr8box}
\begin{keepext}
\exttask A sugar cube is 1\,cm on each side. Calculate its surface area. If it is cut into eight cubes 0.5\,cm on each side, calculate the total surface area. What does this show about cutting a solid into smaller pieces?
\workbox{5cm}
\end{keepext}
\end{yr8box}

% ─── SKILL CHECK ─────────────────────────────────────────────────────────────
\exppart{Skill check}{NavyBlue}

\begin{skillcheck}
\scrow{QP1}{I can write a hypothesis: If the [IV] is \dots, the [DV] will \dots}
\scrow{PC2}{I can name the independent, dependent and controlled variables.}
\scrow{PC8}{I timed each trial with a digital stopwatch.}
\scrow{PA1}{I recorded my times in a table with units in the headings.}
\scrow{PA3}{I can calculate an average and spot an outlier.}
\scrow{PA5}{I can describe a trend using numbers.}
\scrow{EV1}{I can say whether my results support my hypothesis.}
\scrow{EV3}{I can name sources of error and how to reduce them.}
\end{skillcheck}

\task Which skill will you focus on in Experiment 3? Explain what you will do differently.
\answerbox{5cm}
